\documentclass[14pt,a4paper]{extarticle}

\usepackage{fontspec}
\usepackage{polyglossia}
\setdefaultlanguage{russian}
\setotherlanguage{english}

\setmainfont{Tinos}
\setsansfont{Arimo}
\setmonofont{Cousine}[Scale=0.82]

\usepackage{geometry}
\geometry{a4paper, top=2cm, bottom=2cm, left=3cm, right=1.5cm}

\usepackage{amsmath}
\usepackage{amssymb}
\usepackage{graphicx}
\usepackage{float}
\usepackage{booktabs}
\usepackage{xcolor}
\usepackage{listings}
\usepackage{caption} 
\usepackage{tikz}
\usetikzlibrary{positioning}
\usepackage{enumitem}
\usepackage{hyperref}

\hypersetup{hidelinks, unicode}

\linespread{1.3}
\emergencystretch=1.5em
\setlength{\parindent}{1.25cm}

\captionsetup[figure]{name=Рисунок, labelsep=endash}
\captionsetup[table]{name=Таблица, labelsep=endash, singlelinecheck=false, justification=raggedright}
\renewcommand{\lstlistingname}{Листинг}

\newcommand{\Student}{А. А. Устинов}
\newcommand{\Group}{М8О-307БВ-24}
\newcommand{\Teacher}{В. Д. Бахарев}
\newcommand{\Year}{2026}

\definecolor{codegray}{gray}{0.45}
\definecolor{codegreen}{rgb}{0.0,0.45,0.2}
\definecolor{codeblue}{rgb}{0.1,0.2,0.65}

\lstdefinestyle{code}{
	basicstyle=\ttfamily\footnotesize,
	keywordstyle=\color{codeblue}\bfseries,
	commentstyle=\color{codegreen}\itshape,
	stringstyle=\color{codegray},
	numbers=left,
	numberstyle=\tiny\color{codegray},
	numbersep=8pt,
	breaklines=true,
	breakatwhitespace=false,
	showstringspaces=false,
	tabsize=4,
	extendedchars=true,
	inputencoding=utf8,
	frame=single,
	rulecolor=\color{codegray},
	framesep=4pt,
	captionpos=b,
	texcl=true,
	literate={\ }{{\ }}1
}
\lstset{style=code, language=C++}

\lstdefinelanguage{GLSL}[]{C++}{
	morekeywords={vec2,vec3,vec4,mat2,mat3,mat4,layout,location,uniform,in,out,
		push_constant,gl_Position,mix,version,set,binding,std430,readonly,buffer,
		uint,struct,normalize,dot,max,pow,clamp,sqrt,transpose,inverse},
}

\newcommand{\screenshot}[2]{%
	\begin{figure}[H]
		\centering
		\IfFileExists{#1}{%
			\includegraphics[width=0.92\textwidth]{#1}%
		}{%
			\fbox{\parbox[c][45mm][c]{0.9\textwidth}{\centering\itshape
				Место для скриншота\\[2pt]\ttfamily\small #1}}%
		}
		\caption{#2}
	\end{figure}}

\begin{document}

\begin{titlepage}
	\centering

	МОСКОВСКИЙ АВИАЦИОННЫЙ ИНСТИТУТ\par
	(НАЦИОНАЛЬНЫЙ ИССЛЕДОВАТЕЛЬСКИЙ УНИВЕРСИТЕТ)\par

	\vspace{6mm}

	Институт №8 <<Компьютерные науки и прикладная математика>>\par
	Кафедра 806 <<Вычислительная математика и программирование>>\par

	\vspace{6cm}

	{\Large Лабораторная работа №2\par}
	\vspace{3mm}
	{\large по курсу <<Компьютерная графика>>\par}
	\vspace{8mm}
	{\Large\bfseries Камера и освещение\par}
	\vspace{3mm}
	{\large Объект: правильный октаэдр\par}

	\vspace{4cm}

	\begin{flushright}
		\begin{minipage}{0.5\textwidth}
			Выполнил: \Student\par
			Группа: \Group\par
			\vspace{6mm}
			Преподаватель: \Teacher\par
		\end{minipage}
	\end{flushright}

	\vfill
	Москва, \Year\par
\end{titlepage}

\section{Условие}

Цель работы — реализовать виртуальную камеру с возможностью её перемещения и
освещение по модели Блинна-Фонга с использованием нескольких типов источников
света.

Объект — правильный октаэдр из первой лабораторной работы. Программа работает
в реальном времени с использованием Vulkan, GLFW и ImGui. Через
пользовательский интерфейс изменяются параметры камеры (позиция, поворот, угол
обзора), источников света (позиция, направление, цвет, интенсивность и
параметры, свойственные типу источника), рассеянного света и материалов
(альбедо, цвет блика, коэффициент блеска). Нормали вершин генерируются
процедурно, модель Блинна-Фонга вычисляется во фрагментном шейдере, для
точечных источников и прожекторов реализовано затухание по закону обратных
квадратов.

Выполнено базовое задание — отрисовка объекта относительно камеры с одним
направленным источником света — и три дополнительных задания:

\begin{enumerate}[nosep, leftmargin=*]
	\item внешние данные разделены на два набора дескрипторов: набор сцены
	(uniform-буфер с общими параметрами и storage-буферы с массивами
	источников) и набор модели (матрица преобразования и параметры материала);
	\item массив точечных источников в storage-буфере с затуханием по закону
	обратных квадратов, добавлением и удалением источников в интерфейсе;
	\item массив прожекторов в storage-буфере с затуханием, мягкими краями
	освещаемой области, добавлением и удалением источников в интерфейсе.
\end{enumerate}

Для наглядности в сцену добавлен пол — плоскость под октаэдром, на которой
видны пятна света от точечных источников и прожекторов.

% ------------------------------------------------------------------
\section{Метод решения}

\subsection{Камера}

Положение камеры задаётся точкой $\mathbf{e}$, а направление взгляда — двумя
углами: $\psi$ — поворот вокруг вертикальной оси (рысканье, \texttt{yaw}) и
$\varphi$ — наклон вверх или вниз (тангаж, \texttt{pitch}). Единичный вектор
направления взгляда получается из этих углов как точка на единичной сфере:

\begin{equation}
	\mathbf{d} = \bigl(\cos\varphi \sin\psi,\ \ \sin\varphi,\ \ -\cos\varphi \cos\psi\bigr).
	\label{eq:forward}
\end{equation}

При $\psi = \varphi = 0$ получается $\mathbf{d} = (0, 0, -1)$: камера смотрит
вдоль отрицательного направления оси $Z$, как принято в правой системе
координат. Положительное $\psi$ поворачивает взгляд вправо, положительное
$\varphi$ — вверх.

Матрица вида строится тем же способом, что и в первой работе, по положению
камеры, точке наблюдения $\mathbf{e} + \mathbf{d}$ и направлению <<вверх>>
$(0, 1, 0)$, а матрица перспективной проекции $P$ — по углу обзора и отношению
сторон окна. Их произведение $P \cdot V$ передаётся в шейдер как одна матрица.

Угол наклона ограничен диапазоном $\pm 89^\circ$. При $\varphi = \pm 90^\circ$
направление взгляда совпало бы с направлением <<вверх>>, их векторное
произведение обратилось бы в ноль, и базис камеры стал бы неопределённым.

\subsection{Нормали}

Для расчёта освещения в каждой вершине помимо позиции хранится нормаль —
единичный вектор, перпендикулярный поверхности. Нормали генерируются
процедурно: для треугольника с вершинами $\mathbf{a}$, $\mathbf{b}$,
$\mathbf{c}$, перечисленными против часовой стрелки при взгляде снаружи,

\begin{equation}
	\mathbf{n} = \frac{(\mathbf{b} - \mathbf{a}) \times (\mathbf{c} - \mathbf{a})}
	                  {\bigl\|(\mathbf{b} - \mathbf{a}) \times (\mathbf{c} - \mathbf{a})\bigr\|}.
	\label{eq:normal}
\end{equation}

Векторное произведение двух сторон треугольника перпендикулярно его плоскости,
а порядок обхода вершин определяет, что нормаль направлена наружу.

В первой работе октаэдр хранился шестью вершинами: каждая вершина использовалась
четырьмя гранями. Для освещения так нельзя: в каждой вершине октаэдра сходятся
четыре грани с разными нормалями, а атрибут вершины может хранить только одну.
Если усреднить их, освещённость плавно перетекала бы с грани на грань, и
октаэдр выглядел бы скруглённым. Поэтому каждая грань получает свои три
вершины с одинаковой нормалью: 24 вершины для октаэдра и 6 для пола. Такие
нормали называются плоскими.

При переводе в мировые координаты нормаль умножается не на матрицу модели $M$,
а на обратную транспонированную к её верхнему левому блоку $3\times3$:

\begin{equation}
	\mathbf{n}_{\text{мир}} = \bigl(M_{3\times3}^{-1}\bigr)^{\mathsf{T}} \mathbf{n}.
\end{equation}

При неравномерном масштабировании обычная матрица модели искажает нормаль, и
она перестаёт быть перпендикулярной поверхности; обратная транспонированная
матрица этот эффект компенсирует. В данной работе модели только переносятся, и
результат совпал бы, но формула записана в общем виде.

\subsection{Модель освещения Блинна-Фонга}

Обозначения, используемые в формулах освещения, сведены в
таблицу~\ref{tab:light}.

\begin{table}[!htb]
	\caption{Обозначения в формулах освещения}
	\label{tab:light}
	\centering
	\vspace{2mm}
	\begin{tabular}{@{}l p{0.72\textwidth}@{}}
		\toprule
		Обозначение & Что обозначает \\
		\midrule
		$\mathbf{n}$ & единичная нормаль к поверхности в освещаемой точке \\
		$\mathbf{v}$ & единичный вектор от точки к камере \\
		$\mathbf{l}_i$ & единичный вектор от точки к $i$-му источнику \\
		$\mathbf{h}_i$ & полувектор — направление посередине между $\mathbf{l}_i$ и $\mathbf{v}$ \\
		\addlinespace
		$\mathbf{k}_d$ & альбедо — цвет рассеянного (диффузного) отражения материала \\
		$\mathbf{k}_s$ & цвет блика материала \\
		$p$ & коэффициент блеска: чем больше, тем меньше и резче блик \\
		\addlinespace
		$\mathbf{A}$ & рассеянный (фоновый) свет: цвет, умноженный на интенсивность \\
		$\mathbf{E}_i$ & свет, приходящий от $i$-го источника в точку \\
		$\mathbf{C}_i,\ I_i$ & цвет и интенсивность $i$-го источника \\
		$r_i$ & расстояние от точки до $i$-го источника \\
		\addlinespace
		$\odot$ & покомпонентное умножение: каналы R, G и B перемножаются по отдельности \\
		\bottomrule
	\end{tabular}
\end{table}

Цвет точки складывается из рассеянного света и вкладов всех источников:

\begin{equation}
	\mathbf{c} = \mathbf{A} \odot \mathbf{k}_d
	+ \sum_i \Bigl(
		\mathbf{k}_d \max(\mathbf{n} \cdot \mathbf{l}_i,\ 0)
		+ \mathbf{k}_s \max(\mathbf{n} \cdot \mathbf{h}_i,\ 0)^{p}
	\Bigr) \odot \mathbf{E}_i,
	\label{eq:blinn}
\end{equation}

где полувектор

\begin{equation}
	\mathbf{h}_i = \frac{\mathbf{l}_i + \mathbf{v}}{\|\mathbf{l}_i + \mathbf{v}\|}.
	\label{eq:half}
\end{equation}

Первое слагаемое в скобках — диффузное отражение, закон Ламберта: освещённость
пропорциональна косинусу угла между нормалью и направлением на источник.
Поверхность, обращённая к источнику, освещена сильнее всего, повёрнутая боком —
не освещена вовсе. Второе слагаемое — блик: он ярок, когда полувектор
$\mathbf{h}_i$ близок к нормали, то есть когда камера находится рядом с
направлением зеркального отражения луча. Показатель $p$ определяет, насколько
быстро блик гаснет при отклонении от этого направления.

Отличие от исходной модели Фонга — в способе оценки блика. В модели Фонга
вычисляется отражённый луч и его угол с направлением на камеру. Полувектор
вычисляется проще и ведёт себя лучше при скользящих углах падения: угол между
отражённым лучом и направлением на камеру может превысить $90^\circ$, и блик
резко обрывается, тогда как угол между $\mathbf{h}_i$ и нормалью остаётся в
допустимых пределах.

Если поверхность повёрнута от источника ($\mathbf{n} \cdot \mathbf{l}_i \le 0$),
его вклад, включая блик, принимается равным нулю: иначе на неосвещённой стороне
появлялись бы паразитные блики.

\subsection{Типы источников света}

Типы источников различаются только тем, как вычисляются направление
$\mathbf{l}_i$ и приходящий свет $\mathbf{E}_i$.

\textbf{Направленный источник} моделирует очень удалённый источник вроде
солнца: все лучи параллельны, направление $\mathbf{l}$ одинаково для всех точек
сцены, затухания нет:

\begin{equation}
	\mathbf{E} = \mathbf{C} \cdot I.
\end{equation}

\textbf{Точечный источник} излучает одинаково во все стороны из одной точки.
Направление $\mathbf{l}_i$ своё для каждой точки поверхности, а освещённость
убывает по закону обратных квадратов:

\begin{equation}
	\mathbf{E}_i = \frac{\mathbf{C}_i \cdot I_i}{r_i^2}.
	\label{eq:point}
\end{equation}

Физический смысл закона: один и тот же световой поток на расстоянии $r$
распределяется по поверхности сферы площадью $4\pi r^2$, поэтому на единицу
площади приходится доля, обратно пропорциональная $r^2$. При $r \to 0$ формула
даёт бесконечность, поэтому в шейдере квадрат расстояния ограничен снизу
малым числом.

\textbf{Прожектор} — точечный источник, свет которого ограничен конусом с осью
вдоль направления $\mathbf{s}_i$. Для мягкого края задаются два угла,
отсчитываемых от оси: внутренний $\theta_{\text{in}}$, внутри которого яркость
полная, и внешний $\theta_{\text{out}}$, за которым свет пропадает. Между ними
яркость плавно спадает:

\begin{equation}
	\mathbf{E}_i = \frac{\mathbf{C}_i \cdot I_i}{r_i^2} \cdot k_{\text{край}},
	\qquad
	k_{\text{край}} = \operatorname{clamp}\!\left(
		\frac{\cos\theta - \cos\theta_{\text{out}}}
		     {\cos\theta_{\text{in}} - \cos\theta_{\text{out}}},\ 0,\ 1
	\right),
	\label{eq:spot}
\end{equation}

где $\theta$ — угол между осью прожектора и направлением от прожектора на
освещаемую точку, $\cos\theta = -\mathbf{l}_i \cdot \mathbf{s}_i$, а функция
$\operatorname{clamp}(x, 0, 1)$ ограничивает значение отрезком $[0, 1]$.

Внутри внутреннего конуса числитель больше знаменателя и множитель равен
единице; за внешним конусом числитель отрицателен и множитель равен нулю; между
конусами он линейно меняется от единицы до нуля. Сравнение ведётся по
косинусам, а не по самим углам: на отрезке от $0^\circ$ до $90^\circ$ косинус
монотонно убывает, поэтому сравнение косинусов эквивалентно сравнению углов, а
вычислять арккосинус для каждого пикселя не требуется. Косинусы обоих углов
вычисляются один раз на процессоре.

Если внутренний и внешний углы равны, знаменатель обращается в ноль; в шейдере
он ограничен снизу малым числом, и край пятна становится резким.

\subsection{Используемые объекты Vulkan}

В отличие от первой работы, где данные передавались в шейдер через
push-константы, здесь используются наборы дескрипторов. Внешние данные
разделены на два набора (дополнительное задание~1).

\textbf{Набор 0 — данные сцены}, общие для всех объектов:

\begin{itemize}[leftmargin=*, nosep]
	\item привязка 0 — uniform-буфер: матрица $P \cdot V$, положение камеры,
	рассеянный свет, параметры направленного источника и количество
	точечных источников и прожекторов;
	\item привязка 1 — storage-буфер с массивом точечных источников;
	\item привязка 2 — storage-буфер с массивом прожекторов.
\end{itemize}

\textbf{Набор 1 — данные модели}: uniform-буфер с матрицей преобразования и
параметрами материала. У каждой модели свой буфер и свой набор.

При отрисовке набор сцены привязывается один раз за кадр, а набор модели —
перед каждым объектом. Если объектов много, меняется только небольшой набор с
данными модели, а данные сцены остаются привязанными.

Создаются следующие объекты Vulkan:

\begin{itemize}[leftmargin=*]
	\item \texttt{VkDescriptorSetLayout} — две раскладки наборов, описывающие,
	какие ресурсы на каких привязках ожидают шейдеры;

	\item \texttt{VkDescriptorPool} — пул, из которого выделяются наборы: один
	набор сцены и по одному на модель;

	\item \texttt{VkDescriptorSet} — сами наборы, связанные с буферами один раз
	при запуске, поскольку буферы не пересоздаются;

	\item \texttt{VkBuffer} — uniform-буфер сцены, два storage-буфера с
	источниками, а для каждой модели — вершинный буфер и uniform-буфер. Все
	буферы размещены в памяти, доступной процессору, и постоянно отображены в
	его адресное пространство, поэтому данные переписываются каждый кадр
	простым копированием;

	\item \texttt{VkPipelineLayout} — раскладка конвейера с двумя наборами
	дескрипторов;

	\item \texttt{VkPipeline} — графический конвейер. Формат вершины — позиция
	и нормаль, по три числа одинарной точности; остальные настройки (тест
	глубины, отсечение нелицевых граней) совпадают с первой работой.
\end{itemize}

\textbf{Storage-буферы} используются для массивов источников, потому что
только в них массив может быть объявлен в шейдере без указания размера:

\begin{lstlisting}[language=GLSL, numbers=none]
layout(std430, set = 0, binding = 1) readonly buffer PointLights {
	PointLight point_lights[];
};
\end{lstlisting}

В uniform-буфере размер массива обязан быть константой, известной при
компиляции шейдера. Количество действующих источников передаётся в
uniform-буфере сцены, и цикл в шейдере проходит ровно столько элементов. На
стороне приложения буфер выделяется один раз с запасом на 64 источника каждого
типа; это ограничение размера буфера, а не шейдера, и меняется изменением одной
константы.

\textbf{Раскладка данных.} Расположение полей в буферах должно в точности
совпадать с тем, как их читает шейдер. Для uniform-буферов действуют правила
выравнивания \texttt{std140}, для storage-буферов — \texttt{std430}; при
обоих правилах трёхкомпонентный вектор занимает место четырёхкомпонентного.
Поэтому все векторы хранятся как \texttt{vec4}, а свободный четвёртый
компонент используется для смежной величины: интенсивности у цвета источника,
коэффициента блеска у цвета блика.

\textbf{Синхронизация.} Буферы переписываются при записи команд кадра, а не при
обработке интерфейса. К этому моменту стартовый код дожидается завершения
предыдущего кадра, поэтому видеокарта гарантированно не читает буферы в момент
записи.

\subsection{Интерфейс}

Окно интерфейса разделено на разделы: камера, рассеянный свет, направленный
источник, точечные источники, прожекторы и материалы. В разделах источников
каждый элемент списка имеет собственные параметры и кнопку удаления, а под
списком находится кнопка добавления нового источника (дополнительные
задания~2 и~3). Для прожектора внешний угол не может быть сделан меньше
внутреннего. В разделе материалов для каждой модели настраиваются альбедо,
цвет блика и коэффициент блеска.

Ниже приведён фрагментный шейдер, в котором реализована модель освещения.

\lstinputlisting[language=GLSL, caption={Фрагментный шейдер: модель
Блинна-Фонга для трёх типов источников}]{shaders/lit.frag}

% ------------------------------------------------------------------
\section{Результаты}

\screenshot{images/scene.png}{Сцена при запуске: направленный источник,
точечный источник оранжевого цвета и голубой прожектор. Верхние грани октаэдра
освещены сверху, нижние темнее; на полу видны пятна от точечного источника и
прожектора}

\screenshot{images/point_lights.png}{Три точечных источника разных цветов при
выключенном направленном свете. Пятна на полу гаснут с расстоянием по закону
обратных квадратов, грани октаэдра окрашены ближайшими к ним источниками}

\screenshot{images/spot_lights.png}{Два прожектора с разными углами. Голубой
с внутренним углом $10^\circ$ и внешним $30^\circ$ даёт широкий размытый край,
красный с углами $20^\circ$ и $22^\circ$ — почти резкий}

\screenshot{images/material.png}{Изменение материала: зелёное альбедо и
коэффициент блеска 8. При малом коэффициенте блик широкий и растекается по
грани; цвет блика задаётся отдельно от альбедо и остаётся светлым}

% ------------------------------------------------------------------
\section{Выводы}

В ходе работы я научился управлять виртуальной камерой через положение и два
угла поворота, рассчитывать освещение по модели Блинна-Фонга для трёх типов
источников — направленного, точечного и прожектора — и передавать данные в
шейдеры через наборы дескрипторов, uniform- и storage-буферы.

Главное, что я понял: освещение требует от геометрии больше, чем просто
позиции вершин. Нормаль — свойство поверхности, а не вершины, и у октаэдра в
одной вершине сходятся четыре грани с разными нормалями. Поэтому вершины,
которые в первой работе были общими для граней, пришлось продублировать — по
три на каждую грань, — иначе плоские грани освещались бы как скруглённые.

Типы источников различаются лишь тем, как вычисляются направление на источник
и приходящий от него свет, а сама модель отражения для них общая. Мягкий край
прожектора получается линейным переходом между косинусами двух углов, а
сравнение по косинусам позволяет обойтись без вычисления углов в каждом
пикселе.

Разделение данных на набор сцены и набор модели оказалось удобным даже при двух
объектах: данные сцены привязываются один раз за кадр, а для каждого объекта
меняется только небольшой набор с его матрицей и материалом. Storage-буферы
позволили объявить массивы источников в шейдере без фиксированного размера:
количество источников меняется в интерфейсе без изменения шейдера.

\end{document}
